% GENERATED FILE. Edit canonical lessons, then run npm run generate:books.
% canonical-source-hash: fnv1a64:130e90422fb9f341
% canonical-lessons: GU-C463-saral, GU-C463-ochabolu, GU-C463-vacal, GU-C463-bakvu, GU-C463-hasmukhu

\chapter{Chatty and Quiet}
\label{ch:gu-chatty-and-quiet}
\hlchaptermodality{463}

\begin{chapteropening}
\textbf{By the end of this chapter:} \emph{I can say whether someone talks a lot or a little.}

Complete the last lesson of chapter 463: I can say whether someone talks a lot or a little.
\end{chapteropening}

\section[\texorpdfstring{\gu{સરળ} (saraḷ)}{સરળ (saraḷ)}]{\gu{સરળ} (saraḷ) — straightforward (of a person)}

\label{lesson:GU-C463-saral}



\begin{quote}
Before the new one: say the Gujarati for fun-loving, then the Gujarati for to go astray.
\end{quote}

\subsection*{You'll want to know: \gu{સરળ}}
\textbf{\gu{સરળ}} — \emph{saraḷ} — \textquotedblleft{}straightforward (of a person)\textquotedblright{}.

\begin{quote}
\gu{મારી} \gu{બહેનપણી} \gu{સરળ} \gu{સ્વભાવની} \gu{છે}.
\end{quote}

\begin{quote}
\emph{mārī bahenpaṇī saraḷ svabhāvnī che.}
\end{quote}

\begin{quote}
My friend has a straightforward nature.
\end{quote}

\begin{cousinweb}
\textbf{\gu{સરળ}} (\emph{saraḷ}) is a tatsama, Sanskrit \emph{sarala}, \textquotedblleft{}straight, not crooked\textquotedblright{}, as Monier-Williams records it.
\end{cousinweb}

\begin{grammarlens}[title={What you've built}]
One more word for describing what you lived and felt.
\end{grammarlens}

\subsection*{Your turn}
\begin{itemize}
\raggedright
  \item \emph{Say it:} \emph{saraḷ}
  \item \emph{Say it:} \emph{saraḷ}, once more
  \item \emph{Say it:} \emph{mārī bahenpaṇī saraḷ svabhāvnī che.} — My friend has a straightforward nature.
  \item \emph{Recall:} say the Gujarati for fun-loving, then the Gujarati for to go astray, then say \emph{saraḷ} again
\end{itemize}

\begin{tcolorbox}[breakable,colback=teal!4,colframe=teal!35!black,title={Before you move on}]
What does \gu{સરળ} mean? (\textquotedblleft{}Straightforward (of a person)\textquotedblright{}.) Say it once more.
\end{tcolorbox}
\section[\texorpdfstring{\gu{ઓછાબોલું} (ochābolũ)}{ઓછાબોલું (ochābolũ)}]{\gu{ઓછાબોલું} (ochābolũ) — quiet; of few words}

\label{lesson:GU-C463-ochabolu}



\begin{quote}
Before the new one: say the Gujarati for to go astray, then the Gujarati for straightforward.
\end{quote}

\subsection*{You'll want to know: \gu{ઓછાબોલું}}
\textbf{\gu{ઓછાબોલું}} — \emph{ochābolũ} — \textquotedblleft{}quiet; of few words\textquotedblright{}.

\begin{quote}
\gu{અરુણ} \gu{ઓછાબોલો} \gu{છે}, \gu{પણ} \gu{ઘણું} \gu{સમજે} \gu{છે}.
\end{quote}

\begin{quote}
\emph{Aruṇ ochābolo che, paṇ ghaṇũ samje che.}
\end{quote}

\begin{quote}
Arun is a man of few words, but he understands a lot.
\end{quote}

\begin{cousinweb}
\textbf{\gu{ઓછાબોલું}} (\emph{ochābolũ}) is built from \textbf{\gu{ઓછું}} (\emph{ochũ}), \textquotedblleft{}less\textquotedblright{}, and the stem of \textbf{\gu{બોલવું}} (\emph{bolvũ}), \textquotedblleft{}to speak\textquotedblright{}: someone who speaks little. CDIAL lists Gujarati \emph{ochũ} under \emph{*ōccha-} (CDIAL 2540) and \emph{bolvũ} under \emph{*bōll-} (CDIAL 9321).
\end{cousinweb}

\begin{grammarlens}[title={What you've built}]
One more word for describing what you lived and felt.
\end{grammarlens}

\subsection*{Your turn}
\begin{itemize}
\raggedright
  \item \emph{Say it:} \emph{ochābolũ}
  \item \emph{Say it:} \emph{ochābolũ}, once more
  \item \emph{Say it:} \emph{Aruṇ ochābolo che, paṇ ghaṇũ samje che.} — Arun is a man of few words, but he understands a lot.
  \item \emph{Recall:} say the Gujarati for to go astray, then the Gujarati for straightforward, then say \emph{ochābolũ} again
\end{itemize}

\begin{tcolorbox}[breakable,colback=teal!4,colframe=teal!35!black,title={Before you move on}]
What does \gu{ઓછાબોલું} mean? (\textquotedblleft{}Quiet; of few words\textquotedblright{}.) Say it once more.
\end{tcolorbox}
\section[\texorpdfstring{\gu{વાચાળ} (vācāḷ)}{વાચાળ (vācāḷ)}]{\gu{વાચાળ} (vācāḷ) — talkative}

\label{lesson:GU-C463-vacal}



\begin{quote}
Before the new one: say the Gujarati for straightforward, then the Gujarati for of few words.
\end{quote}

\subsection*{You'll want to know: \gu{વાચાળ}}
\textbf{\gu{વાચાળ}} — \emph{vācāḷ} — \textquotedblleft{}talkative\textquotedblright{}.

\begin{quote}
\gu{વાચાળ} \gu{મામી} \gu{આખો} \gu{દિવસ} \gu{બોલે} \gu{છે}.
\end{quote}

\begin{quote}
\emph{vācāḷ māmī ākho divas bole che.}
\end{quote}

\begin{quote}
My talkative aunt talks all day long.
\end{quote}

\begin{cousinweb}
\textbf{\gu{વાચાળ}} (\emph{vācāḷ}) is a tatsama, Sanskrit \emph{vācāla}, \textquotedblleft{}talkative, chattering\textquotedblright{}, as Monier-Williams records it.
\end{cousinweb}

\begin{grammarlens}[title={What you've built}]
One more word for describing what you lived and felt.
\end{grammarlens}

\subsection*{Your turn}
\begin{itemize}
\raggedright
  \item \emph{Say it:} \emph{vācāḷ}
  \item \emph{Say it:} \emph{vācāḷ}, once more
  \item \emph{Say it:} \emph{vācāḷ māmī ākho divas bole che.} — My talkative aunt talks all day long.
  \item \emph{Recall:} say the Gujarati for straightforward, then the Gujarati for of few words, then say \emph{vācāḷ} again
\end{itemize}

\begin{tcolorbox}[breakable,colback=teal!4,colframe=teal!35!black,title={Before you move on}]
What does \gu{વાચાળ} mean? (\textquotedblleft{}Talkative\textquotedblright{}.) Say it once more.
\end{tcolorbox}
\section[\texorpdfstring{\gu{બકવું} (bakvũ)}{બકવું (bakvũ)}]{\gu{બકવું} (bakvũ) — to babble; to talk nonsense}

\label{lesson:GU-C463-bakvu}



\begin{quote}
Before the new one: say the Gujarati for of few words, then the Gujarati for talkative.
\end{quote}

\subsection*{You'll want to know: \gu{બકવું}}
\textbf{\gu{બકવું}} — \emph{bakvũ} — \textquotedblleft{}to babble; to talk nonsense\textquotedblright{}.

\begin{quote}
\gu{તાવમાં} \gu{દર્દી} \gu{કંઈક} \gu{બકતો} \gu{હતો}.
\end{quote}

\begin{quote}
\emph{tāvmā̃ dardī kaīk bakto hato.}
\end{quote}

\begin{quote}
In his fever the patient was babbling something.
\end{quote}

Past (he …): \gu{તે} \gu{ઘણું} \gu{બક્યો}, \emph{te ghaṇũ bakyo}.

\begin{cousinweb}
\textbf{\gu{બકવું}} (\emph{bakvũ}) goes back to a reconstructed \emph{*bakk-}, \textquotedblleft{}chatter\textquotedblright{}, an imitative word; CDIAL lists the Gujarati verb, \textquotedblleft{}to prattle\textquotedblright{} (CDIAL 9117).
\end{cousinweb}

\begin{grammarlens}[title={What you've built}]
One more word for describing what you lived and felt.
\end{grammarlens}

\subsection*{Your turn}
\begin{itemize}
\raggedright
  \item \emph{Say it:} \emph{bakvũ}
  \item \emph{Say it:} \emph{bakvũ}, once more
  \item \emph{Say it:} \emph{tāvmā̃ dardī kaīk bakto hato.} — In his fever the patient was babbling something.
  \item \emph{Recall:} say the Gujarati for of few words, then the Gujarati for talkative, then say \emph{bakvũ} again
\end{itemize}

\begin{tcolorbox}[breakable,colback=teal!4,colframe=teal!35!black,title={Before you move on}]
What does \gu{બકવું} mean? (\textquotedblleft{}To babble; to talk nonsense\textquotedblright{}.) Say it once more.
\end{tcolorbox}
\section[\texorpdfstring{\gu{હસમુખું} (hasmukhũ)}{હસમુખું (hasmukhũ)}]{\gu{હસમુખું} (hasmukhũ) — cheerful; smiling}

\label{lesson:GU-C463-hasmukhu}



\begin{quote}
Before the new one: say the Gujarati for talkative, then the Gujarati for to babble.
\end{quote}

\subsection*{You'll want to know: \gu{હસમુખું}}
\textbf{\gu{હસમુખું}} — \emph{hasmukhũ} — \textquotedblleft{}cheerful; smiling\textquotedblright{}.

\begin{quote}
\gu{અમારી} \gu{નર્સ} \gu{હંમેશા} \gu{હસમુખી} \gu{રહે} \gu{છે}.
\end{quote}

\begin{quote}
\emph{amārī nars hammeśā hasmukhī rahe che.}
\end{quote}

\begin{quote}
Our nurse is always cheerful.
\end{quote}

\begin{cousinweb}
\textbf{\gu{હસમુખું}} (\emph{hasmukhũ}) is built from the stem of \textbf{\gu{હસવું}} (\emph{hasvũ}), \textquotedblleft{}to laugh\textquotedblright{}, and \emph{mukh}, \textquotedblleft{}face\textquotedblright{}, Sanskrit \emph{mukha} (Monier-Williams): a laughing face.
\end{cousinweb}

\begin{grammarlens}[title={What you've built}]
One more word for describing what you lived and felt.
\end{grammarlens}

\subsection*{Your turn}
\begin{itemize}
\raggedright
  \item \emph{Say it:} \emph{hasmukhũ}
  \item \emph{Say it:} \emph{hasmukhũ}, once more
  \item \emph{Say it:} \emph{amārī nars hammeśā hasmukhī rahe che.} — Our nurse is always cheerful.
  \item \emph{Recall:} say the Gujarati for talkative, then the Gujarati for to babble, then say \emph{hasmukhũ} again
\end{itemize}

\begin{tcolorbox}[breakable,colback=teal!4,colframe=teal!35!black,title={Before you move on}]
What does \gu{હસમુખું} mean? (\textquotedblleft{}Cheerful; smiling\textquotedblright{}.) Say it once more.
\end{tcolorbox}
